\documentclass[11pt,a4paper]{article}

% Compile with XeLaTeX
\usepackage{fontspec}
\usepackage{polyglossia}
\setdefaultlanguage{bulgarian}
\setotherlanguage{english}
\defaultfontfeatures{Ligatures=TeX}
\setmainfont{Times New Roman}
\setsansfont{Arial}
\setmonofont{Menlo}
\newfontfamily\cyrillicfont{Times New Roman}
\newfontfamily\cyrillicfontsf{Arial}
\newfontfamily\cyrillicfonttt{Menlo}
\usepackage[a4paper,margin=2.5cm]{geometry}
\usepackage{graphicx}
\usepackage{float}
\usepackage{array}
\usepackage{booktabs}
\usepackage{tabularx}
\usepackage{setspace}
\usepackage{fancyhdr}
\usepackage{lastpage}
\usepackage{xcolor}
\usepackage{hyperref}
\usepackage{caption}
\usepackage{enumitem}
\usepackage{titlesec}
\usepackage{parskip}

\definecolor{portfolio}{HTML}{2F4858}
\definecolor{portfolioaccent}{HTML}{33D2FF}
\definecolor{portfoliohome}{HTML}{F2E9DE}
\definecolor{portfolioabout}{HTML}{FFC1B4}
\definecolor{portfoliocontacts}{HTML}{92D9E8}
\definecolor{portfolioink}{HTML}{0F1108}

\hypersetup{
    colorlinks=true,
    linkcolor=portfolio,
    urlcolor=portfolio,
    pdftitle={Курсова работа по CSCB857 - Портфолио сайт Dian Kalaydzhiev},
    pdfauthor={Диян Калайджиев}
}

\newcommand{\setassetpath}[1]{%
  \IfFileExists{latex/#1}{%
    \def\assetfile{latex/#1}%
  }{%
    \IfFileExists{#1}{%
      \def\assetfile{#1}%
    }{%
      \def\assetfile{../latex/#1}%
    }%
  }%
}

\captionsetup[figure]{
    name=Фигура,
    font=small,
    labelfont=bf,
    justification=centering
}

\setstretch{1.12}
\setlength{\parindent}{1.25em}
\emergencystretch=1.5em
\setlist[itemize]{leftmargin=1.5em,itemsep=0.3em}

\pagestyle{fancy}
\fancyhf{}
\fancyhead[L]{CSCB857}
\fancyhead[R]{Съвременни тенденции в Web-дизайна}
\fancyfoot[C]{\thepage\ /\ \pageref{LastPage}}
\setlength{\headheight}{15pt}
\renewcommand{\headrulewidth}{0.4pt}
\renewcommand{\footrulewidth}{0pt}

\titleformat{\section}{\Large\bfseries\color{portfolio}}{\thesection.}{0.6em}{}
\titleformat{\subsection}{\large\bfseries\color{portfolioink}}{\thesubsection.}{0.5em}{}
\titleformat{\subsubsection}{\normalsize\bfseries\color{portfolio}}{\thesubsubsection.}{0.5em}{}

\newcommand{\portfoliofigure}[4][0.94\textwidth]{
\begin{figure}[H]
\centering
\setassetpath{#2}
\includegraphics[width=#1,height=0.58\textheight,keepaspectratio]{\assetfile}
\caption{#3}
\label{#4}
\end{figure}
}

\begin{document}

\begin{titlepage}
    \centering
    \vspace*{2cm}
    {\Large Нов български университет\par}
    \vspace{0.8cm}
    {\large Департамент „Информатика“\par}
    \vfill
    {\Huge\bfseries Курсова работа\par}
    \vspace{0.5cm}
    {\Large по CSCB857\\[0.2cm]Съвременни тенденции в Web-дизайна\par}
    \vspace{1.2cm}
    {\LARGE\bfseries Документация и анализ на портфолио сайт\\[0.2cm]\textit{Dian Kalaydzhiev}\par}
    \vspace{1.2cm}
    \begin{tabular}{>{\bfseries}p{4.3cm}p{8cm}}
        Автор: & Диян Калайджиев \\
        Факултетен номер: & f108356 \\
        Курс / семестър: & 4-ти курс / 8-ми семестър \\
        Тип проект: & Едностраничен портфолио уеб сайт \\
        Основни технологии: & HTML, CSS, JavaScript, React \\
        Предназначение: & Представяне на професионален профил, умения и контакти \\
    \end{tabular}
    \vfill
    {\large София\par}
    {\large 29 април 2026 г.\par}
\end{titlepage}

\tableofcontents
\clearpage

\section{Въведение}
Настоящият документ представя цялостен анализ на личен портфолио сайт, разработен като едностранично уеб приложение с четири ясно обособени секции: \textit{Home}, \textit{About}, \textit{Certificates} и \textit{Contacts}. Основната му цел е да демонстрира професионалния профил на автора по ясен, визуално подреден и лесен за възприемане начин.

Сайтът е изграден с HTML, CSS и JavaScript, като React е използван само в началната секция, където интеракцията има най-голяма стойност за общото впечатление. Това решение е важно, защото показва прагматичен подбор на технологии: там, където статичната структура е достатъчна, не се добавя излишна сложност, а там, където се търси по-емоционално и динамично поведение, се използва компонентен подход. По този начин сайтът остава лек като архитектура, но едновременно с това постига по-висока визуална разпознаваемост.

Документацията разглежда не само техническата организация на проекта, но и начина, по който отделните секции работят като части от общо потребителско изживяване. Особен акцент е поставен върху визуалната логика, адаптивността, навигацията и връзката между съдържание и интеракция. Добавените екранни снимки са подредени в реда, в който потребителят реално среща съответните екрани и състояния, така че анализът да следва естествения поток на разглеждане.

\section{Архитектура на проекта и карта на сайта}
Сайтът е реализиран като \textbf{single-page} структура. Вместо да разделя информацията в множество самостоятелни HTML документи, проектът обединява цялото съдържание в една входна страница и използва навигация с котви за бързо придвижване между основните секции. Това решение е особено подходящо за персонално портфолио, защото поддържа вниманието върху автора и избягва прекъсването, което често възниква при зареждане на отделни страници.

Линейната архитектура създава предвидим и лесен за разбиране потребителски поток. Още от първия екран посетителят получава визуално силно въведение, след което преминава към биографична информация, доказателства за квалификация и накрая към конкретни канали за връзка. Тази последователност не е случайна: първо се изгражда интерес, после се подава контекст, след това се демонстрира надеждност и накрая се предлага действие.

\portfoliofigure[0.98\textwidth]{01_diagram.png}{Обща последователност на основните секции в сайта}{fig:sitemap}

Както се вижда на фигура~\ref{fig:sitemap}, структурата е логически подредена и не изисква допълнително обучение от страна на потребителя. Това е силна страна на проекта, защото личното портфолио трябва да бъде възприемано в рамките на секунди. При подобен тип сайт сложната навигация по-скоро би разсеяла посетителя, отколкото да му помогне.

Допълнително предимство на този архитектурен модел е, че визуалната идентичност остава постоянна по време на цялото разглеждане. Цветовите секции се сменят, но сайтът не губи своя ритъм и цялост. Именно този плавен преход между отделните блокове е една от характерните особености на модерния \textit{single-page} подход.

\subsection{Файлова организация}
Кодът е разделен по отговорности. Основният HTML документ дефинира общата структура и зарежда външните зависимости, CSS е разпределен по секции, а JavaScript е отделен в няколко файла според вида поведение. Тази организация е добра инженерна практика, защото позволява всяка част да се редактира сравнително самостоятелно, без да се създава ненужно преплитане между слоевете.

\begin{table}[H]
\centering
\caption{Основни файлове и тяхната роля}
\begin{tabularx}{\textwidth}{>{\raggedright\arraybackslash}p{4.6cm}>{\raggedright\arraybackslash}X}
\toprule
Файл & Роля в проекта \\
\midrule
\path{index.html} & Основен документ, в който са дефинирани секциите, навигацията и CDN зависимостите. \\
\path{styles/style.css} & Централен CSS файл, който импортира модулите за всяка секция. \\
\path{styles/home.css} & Подреждане на началния екран, навигация, изображение и позициониране на текстовите елементи. \\
\path{styles/about.css} & Стилове за секцията с биографична информация и бутона за CV. \\
\path{styles/certificates.css} & Визуално оформление на сертификатите и слайдера. \\
\path{styles/contacts.css} & Карти за социални и контактни канали, както и уведомлението за копиране. \\
\path{styles/mobileMenu.css} & Отделна анимирана навигация за мобилна версия. \\
\path{script.js} & React приложението за началния екран и допълнителни функции като копиране на имейл. \\
\path{js/scrollFunc.js} & Оцветяване на активния навигационен елемент според текущата секция. \\
\path{js/switchText.js} & Typewriter ефект за сменящия се текст в началния екран. \\
\bottomrule
\end{tabularx}
\end{table}

Разделянето на стиловете по файлове е особено полезно за проект с различни визуални секции. \textit{Home} има герой-композиция и интерактивен характер, \textit{About} разчита на текст и илюстрация, \textit{Certificates} работи като витрина, а \textit{Contacts} като функционална решетка. Ако всички тези правила бяха събрани в един голям файл, поддръжката щеше да бъде значително по-трудна и рискът от странични ефекти при промени щеше да е по-висок.

\section{Визуална концепция и дизайнерски решения}
Визуалната посока на сайта се основава на ясно разграничени цветни полета, заоблени форми, достатъчно свободно пространство и отчетлив контраст между секциите. Вместо да разчита на тежки декоративни ефекти, проектът използва чисти фонове и силни цветови блокове, чрез които всяка част от страницата получава собствен характер. Това е актуален подход в съвременния web дизайн, особено при портфолио сайтове, където съдържанието трябва да изглежда уверено и целенасочено.

Избраната палитра има и функционална роля. Топлият светъл фон в началната секция създава достъпно и приветливо първо впечатление. Секцията \textit{About} преминава към по-плътен розово-коралов тон, който визуално задържа вниманието върху текста. \textit{Certificates} използва тъмен фон, чрез който изображенията на сертификатите изпъкват по-силно, а финалната секция \textit{Contacts} се връща към по-лека и открита атмосфера със светлосиня основа. Така цветът не е само декоративен избор, а инструмент за навигация и емоционално темпериране.

Типографският избор с \textit{Montserrat} също е последователен спрямо характера на сайта. Шрифтът работи добре при кратки заглавия, подчертано професионални послания и интерфейсни елементи. Той допринася за по-съвременен и технологично ориентиран образ, без да жертва четимостта. Комбинацията от едри заглавия, по-къси текстови блокове и ясно отделени бутони позволява информацията да бъде сканирана бързо, което е ключово за портфолио, разглеждано от работодатели или клиенти.

Друг важен аспект е модулността на визуалното решение. Централният стилов файл не се опитва да описва целия сайт самостоятелно, а само организира отделните CSS модули. Това подсказва зряло отношение към поддръжката на интерфейса. В резултат на това всяка секция може да бъде доразвивана като отделен визуален компонент, без да се разрушава общата логика на проекта.

\section{Начална секция Home}
Началният екран е най-силно изразената част от проекта и носи основната тежест на първото впечатление. Още при зареждане потребителят вижда едновременно фотографски елемент, кратък текстов блок и интерактивен компонент, който придава характер на иначе стандартна портфолио структура. Това е добро решение, защото началният екран трябва не само да информира, но и да отличи автора сред множество сходни представяния.

\portfoliofigure{02_homepage.png}{Началната секция \textit{Home} с изображение, типографски блок и интерактивен бутон}{fig:home-screen}

Както се вижда на фигура~\ref{fig:home-screen}, композицията е изградена около ясен баланс между лице на автора и кратко, лесно за запомняне послание. Лявата част изгражда персонално присъствие, а дясната част съдържа динамичен текст и бутон, който провокира взаимодействие. Фоновият слой с технологични елементи подсилва професионалната идентичност, без да доминира над основното съдържание.

Тази секция функционира като \textit{hero} зона в съвременния смисъл на думата. Тя не е просто заглавие в началото на страницата, а самостоятелно преживяване, което заявява тон, ниво на увереност и визуална култура. Именно тук присъствието на React е най-оправдано, защото интерактивността е част от посланието, а не само техническо усложнение.

\subsection{HTML реализация}
HTML структурата на началната секция е минималистична. В основата си тя съдържа контейнер \texttt{\#app}, в който React рендерира съдържанието. Подобен подход прави \textit{index.html} по-чист и отделя динамичната визуализация в самостоятелен слой. Това е полезно, защото стартовият документ остава лесен за четене, а интерактивното поведение се концентрира на едно място.

От гледна точка на поддръжката този избор също е удачен. Ако в бъдеще началният екран бъде променен изцяло или заменен с друга интерактивна концепция, останалите секции на сайта могат да останат почти непокътнати. Това означава, че най-експерименталната част на портфолиото е сравнително добре капсулирана.

\subsection{CSS оформление}
CSS позиционира снимката, заглавията и интерактивната зона така, че екранът да изглежда цялостно още при първия поглед. Използвани са абсолютни позиции и стойности, зависими от ширината на екрана, което позволява да се постигне визуална драматичност при по-големи дисплеи. Този тип оформление е характерен за съвременни портфолиа, при които началната секция трябва да има повече сценично присъствие.

Същевременно дизайнът не разчита само на големи елементи. Има ясно разстояние между отделните части, което пази композицията от претрупване. Това е особено важно при силен начален екран, защото ако всяко нещо се стреми да бъде акцент, накрая нищо не остава водещо. Тук приоритетът е добре структуриран: първо лице, после текст, после действие.

\subsection{React логика}
React се използва за бутона, анимацията на ръката и следенето на различни състояния като \textit{hover}, позиция на курсора и временно променящ се текст върху бутона. Това е най-съвременният технически слой в проекта и ясно показва желание интеракцията да бъде описана декларативно, а не чрез множество отделни директни манипулации върху DOM.

\portfoliofigure[0.58\textwidth]{07_animation.png}{Състояние на интерактивната анимация в React компонента на началния екран}{fig:home-react-animation}

Както се вижда на фигура~\ref{fig:home-react-animation}, интерактивният елемент не е декоративно допълнение, а част от начина, по който началната секция комуникира характер и динамика.

Важно е да се отбележи, че React не е използван като самоцел. Ако цялата страница беше пренаписана в React без необходимост, това би увеличило сложността без съществена полза. Тук обаче библиотеката е ограничена до точно този компонент, в който промяната на състоянията има визуална стойност и допринася за характерното усещане на началния екран.

\subsection{JavaScript поведение}
Освен React, в началната секция присъства и отделен \textit{typewriter} ефект, реализиран с vanilla JavaScript. Този детайл добавя движение към текста под основното приветствие и създава чувство, че началният екран е жив, а не статичен. При портфолио сайт това е особено полезно, защото кратките въртящи се послания могат да представят няколко важни професионални роли в ограничено пространство.

Комбинацията между React и самостоятелен JavaScript модул е интересна и от инженерна гледна точка. Тя показва, че проектът не следва догматично една технология навсякъде, а използва различни средства според конкретната нужда. Това е прагматичен подход, който често води до по-икономична и по-разбираема реализация.

\section{Секция About}
Секцията \textit{About} има задачата да прехвърли потребителя от емоционалното първо впечатление към по-съдържателно професионално представяне. Тук фокусът вече не е върху играта с интерфейса, а върху кратки, ясно формулирани твърдения за позиция, опит и академично присъствие. Подобна смяна на тона е логична, защото след първоначалното привличане на вниманието сайтът трябва да предостави основания за доверие.

\portfoliofigure{03_about_page.png}{Секция \textit{About} с кратко професионално представяне и бутон за CV}{fig:about-screen}

На фигура~\ref{fig:about-screen} се вижда композиция, която използва достатъчно празно пространство и ясно разположени елементи. В лявата част е събран текстовият блок, а вдясно стои илюстративният елемент. Това решение работи добре, защото текстът не е затрупан от графика, но и секцията не остава прекалено суха или документална.

Силната страна на тази част е, че съдържанието е организирано като бързо сканиращ се списък от професионални маркери. Вместо дълъг автобиографичен абзац посетителят получава отделни твърдения, които могат да бъдат възприети почти мигновено. За портфолио сайт това е по-ефективен подход, понеже повечето потребители първо сканират, а едва след това задълбават.

\subsection{HTML структура}
HTML на секцията е прост и функционален. Текстовите редове са групирани заедно, а бутонът за изтегляне на CV стои като самостоятелен призив за действие. Илюстрацията е отделен елемент, което прави композицията ясна както визуално, така и на ниво структура. Подобно разделение между информация и визуална опора е лесно за поддръжка и позволява по-гъвкави промени в бъдеще.

Наличието на директен бутон за CV е особено силен UX избор. Потребителят не трябва да търси допълнителни връзки или менюта, за да стигне до по-подробния професионален документ. Действието е поставено точно там, където възниква естествен интерес към биографията.

\subsection{CSS и адаптивност}
CSS използва \texttt{flex} оформление за хоризонтално подреждане на текста и изображението, а при по-малки екрани преминава към по-опростен вертикален или блоков режим. Това е правилна стратегия, защото секцията трябва да пази четимостта на съдържанието дори когато няма достатъчно хоризонтално пространство за изразена двуколонна композиция.

Визуалният акцент върху бутона за CV също е премерен. Той е достатъчно видим, за да бъде разпознат като действие, но не засенчва основния текст. Така сайтът поддържа усещане за професионална увереност, а не за агресивна самореклама. Този баланс е важен за добро личностно представяне в дигитална среда.

\section{Секция Certificates}
Секцията \textit{Certificates} играе ролята на визуално доказателство за натрупани знания и практически обучения. Тя не разчита на дълъг текст, а на показване на реални сертификати в подвижен слайдер. Това е подход, който е особено ефективен за портфолио сайтове, защото превръща квалификациите в лесно възприемаема галерия, вместо в сух списък от курсове и години.

\portfoliofigure{04_certifiates.png}{Секция \textit{Certificates} с автоматично движещ се слайдер от сертификати}{fig:certificates-screen}

Фигура~\ref{fig:certificates-screen} показва тъмна секция, в която сертификатите изпъкват силно върху фона. Контрастът е добре подбран: изображенията придобиват по-голяма тежест и се възприемат като самостоятелни артефакти на професионално развитие. Така секцията постига усещане за представителност, което трудно би се получило при обикновен списък от текстови връзки.

Друг важен плюс е, че големият брой сертификати не претоварва потребителя наведнъж. Слайдерът въвежда ритъм и дозиране на информацията. Посетителят разбира, че има значителен набор от постижения, без да се сблъсква с дълга вертикална стена от изображения.

\subsection{HTML структура}
HTML реализацията следва логиката на библиотеката Splide с ясен външен контейнер, \texttt{track} област и списък от \texttt{slide} елементи. Това е добра практика, защото структурата остава четима и съвместима с вече утвърден компонентен модел. Всеки сертификат е отделен визуален елемент, което прави съдържанието лесно за разширяване при добавяне на нови постижения.

Фактът, че изображенията са представени като поредица от еднотипни елементи, улеснява и поддръжката. Ако в бъдеще се наложи добавяне, премахване или пренареждане на сертификати, промяната ще бъде локална и предвидима. Това е дребен, но важен белег за добра вътрешна организация.

\subsection{JavaScript конфигурация}
Слайдерът се конфигурира чрез JavaScript с циклично поведение, свободно плъзгане и автоматично скролиране. Тази комбинация е добре избрана за съдържание от тип галерия, защото създава усещане за непрекъснатост и не изисква потребителят непременно да взаимодейства активно, за да види богатството на секцията.

Допълнително конфигурацията променя броя видими елементи според ширината на екрана. Това е съществен детайл за адаптивността. Без подобна логика сертификатите биха изглеждали или твърде дребни на малък екран, или твърде разтегнати на голям. Тук поведението е съобразено с реалната среда на употреба.

\subsection{CSS оформление}
Стиловете на секцията определят не само размерите и разстоянията между картините, но и цялостния ритъм на възприемане. Големите визуални полета, равномерните отстояния и тъмният фон превръщат секцията в своеобразна витрина. Вместо да изглежда като техническо приложение на библиотека, тя се възприема като част от визуалната идентичност на сайта.

Този ефект е важен и от комуникационна гледна точка. Сертификатите не са второстепенна информация, а част от доверието, което сайтът изгражда. Затова начинът, по който са показани, трябва да изглежда добре обмислен и представителен, а не случаен.

\section{Секция Contacts}
Секцията \textit{Contacts} затваря потребителския поток и превръща натрупаното впечатление в практическа възможност за действие. След като посетителят е преминал през личното представяне, биографичната информация и сертификатите, тук той трябва максимално бързо да стигне до удобен канал за връзка. Именно затова секцията е визуално опростена и ориентирана към ясно разпознаваеми икони.

\portfoliofigure{05_contacts.png}{Секция \textit{Contacts} с подредени канали за връзка и уведомление за копиране}{fig:contacts-screen}

На фигура~\ref{fig:contacts-screen} се вижда светъл фон и равномерно разположени контактни елементи. След тъмната и визуално по-тежка секция със сертификати това решение действа освежаващо и естествено насочва вниманието към действията. Финалният екран не се опитва да впечатлява със сложност, а да завърши пътя чисто и безпроблемно.

Силна страна тук е, че контактите са представени като визуални карти, а не като дълги текстови редове. Това намалява когнитивното натоварване и прави секцията лесна за използване както на десктоп, така и на мобилни устройства. Посетителят вижда познати символи и реагира почти инстинктивно.

\subsection{HTML структура}
HTML съдържа отделни \texttt{div} блокове с изображения на социални платформи и имейл. Някои от тях отварят външни профили, а един елемент активира функция за копиране на имейл адреса. Това е директен и разбираем подход, който не изисква допълнителни слоеве на навигация.

Наличието на отделен контейнер за уведомление е добър детайл. Вместо потребителят да предполага, че действието е успяло, интерфейсът връща моментално потвърждение. Това прави взаимодействието по-надеждно и намалява вероятността от колебание дали имейлът действително е копиран.

\subsection{CSS оформление}
CSS използва \texttt{flex-wrap}, за да подрежда иконите в последователни редове и да позволява естествено пренареждане при по-малки ширини. Изображенията са стилизирани с еднакъв визуален ритъм, което придава завършен вид на секцията, въпреки че всяка платформа има собствена иконография.

Важно е, че визуалното решение не пренатоварва финалния екран. Контактната секция трябва да бъде бърза и удобна, а не декоративно сложна. Тук формата е подчинена на функцията, но без да се губи цялостният стил на сайта.

\subsection{JavaScript интеракция}
Най-съществената логика в тази секция е свързана с копиране на имейл към клипборда и кратко уведомление след успешното действие. Макар функционалността да е малка по обем, тя има висок практически ефект. Потребителят може веднага да използва адреса, без да го маркира и преписва ръчно.

Това е добър пример как сравнително дребна интеракция подобрява усещането за завършеност. При личен сайт точно такива на пръв поглед малки удобства често правят разликата между просто красива страница и действително използваем интерфейс.

\section{Навигация, скрол логика и мобилно меню}
Навигацията е ключова за едностраничен сайт, защото тук тя не просто прехвърля към нов адрес, а служи като карта на цялото съдържание. В проекта са реализирани два паралелни режима: стандартна фиксирана десктоп навигация и самостоятелна мобилна навигация с анимирано отваряне. Това позволява интерфейсът да остане удобен в различни контексти на употреба, без да жертва визуалната си идентичност.

\portfoliofigure[0.48\textwidth]{06_mobile_nav.png}{Мобилна навигация с анимирано разгъване и ясно подредени връзки}{fig:navigation-screen}

Фигура~\ref{fig:navigation-screen} показва мобилния режим като отделен интерфейсен жест, а не просто свит вариант на десктоп менюто. Линковете са достатъчно големи за работа с докосване, а анимацията създава отчетливо усещане за преход между затворено и отворено състояние. Това е важен детайл, защото в мобилна среда навигацията трябва да бъде едновременно видима и ненатрапчива.

При десктоп изгледа, от друга страна, менюто остава по-дискретно и служи като постоянна хоризонтална ориентирна линия. Този контраст между двата режима е признак за добре обмислена адаптация, а не просто за механично скриване и показване на едни и същи елементи.

\subsection{HTML структура}
В \texttt{index.html} присъстват две отделни навигации: \texttt{\#desktopNav} и \texttt{\#mobileNav}. Това решение е директно и практично. Вместо JavaScript да създава и унищожава елементи според ширината на екрана, видимостта се управлява главно чрез CSS. Така структурата остава предвидима и сравнително лесна за дебъгване.

Подобен подход е удачен за сравнително малък проект, защото намалява сложността на зависимостите между layout и поведение. В същото време оставя достатъчно свобода двата режима да бъдат визуално различни, когато това е необходимо.

\subsection{JavaScript при скрол}
Скрол логиката следи активната секция и добавя клас към съответния елемент от навигацията. Това е важно за ориентацията в дълга едностранична структура, защото потребителят получава постоянна обратна връзка къде точно се намира. Подобно поведение е типично за качествени \textit{single-page} интерфейси и създава усещане за по-стабилна информационна архитектура.

Практическата полза е особено голяма при преминаване между визуално различни секции. Когато фонът и композицията се сменят рязко, активният елемент в менюто играе ролята на втори навигационен маркер и намалява риска посетителят да загуби ориентир.

\subsection{CSS анимация на мобилното меню}
Мобилното меню използва CSS-базирана техника с checkbox, чрез която състоянието на навигацията се контролира почти изцяло без допълнителна JavaScript логика. Това е елегантно решение, защото запазва интеракцията надеждна и лесна за поддръжка. Анимацията при отваряне не е просто визуален ефект, а начин да се направи преходът между състоянията по-ясен за потребителя.

Формата на разгъването също не е случайна. Вместо стандартно плъзгащ се панел, менюто използва по-органичен визуален жест, който допринася за отличителния характер на сайта. Това помага мобилният режим да бъде възприет като пълноценна част от дизайна, а не като второстепенно техническо допълнение.

\section{Адаптивност, достъпност и инженерни изводи}
Проектът използва последователно \textit{media queries} в отделните CSS файлове. Това показва, че адаптивността е мислена не като последен етап на корекция, а като естествено изискване при изграждането на всяка секция. Този подход е правилен, защото \textit{Home}, \textit{About}, \textit{Certificates} и \textit{Contacts} имат различна композиционна логика и не могат да бъдат пренаредени надеждно чрез един универсален шаблон.

Особено добро решение е, че адаптивността не се свежда единствено до оразмеряване на елементите. При по-малки екрани реално се променя начинът, по който секциите са подредени и възприемани. Например \textit{About} преминава към по-опростена структура, а сертификатният слайдер редуцира броя едновременно видими елементи. Така интерфейсът не просто се свива, а се адаптира смислено.

От гледна точка на достъпността могат да се отбележат няколко положителни страни. Използвани са \texttt{alt} атрибути за изображенията, навигационните връзки са ясно именувани, а основните секции са логически разделени. Това създава добра базова структура и улеснява както визуалното възприемане, така и машинното разбиране на документа от помощни технологии.

Съществуват и посоки за бъдещо подобрение. Част от интерактивните области могат да бъдат изведени чрез по-семантични бутони и връзки вместо \texttt{onclick} върху контейнерни елементи. Също така сайтът би спечелил от по-видими \texttt{:focus} състояния за потребители, които навигират с клавиатура. Това не отменя качествата на текущата реализация, а по-скоро очертава следващото ниво на зрялост на проекта.

От инженерна гледна точка най-силните решения могат да бъдат обобщени така:
\begin{itemize}
    \item ясна модулност на стиловете и разделяне на файловете по секции;
    \item използване на React само там, където интерактивността наистина носи стойност;
    \item прагматично включване на външна библиотека за слайдер вместо ненужно преизобретяване;
    \item добра визуална сегментация на съдържанието чрез цвят, композиция и ритъм;
    \item балансирано съчетаване на статични и динамични елементи в рамките на един сравнително лек проект.
\end{itemize}

Към това може да се добави и още един важен извод: проектът успява да остане личен и разпознаваем, без да губи функционалност. Това е съществен критерий за добро портфолио. Не е достатъчно сайтът само да изглежда модерно; той трябва да помага на потребителя бързо да разбере кой е авторът, какъв опит има и как може да бъде потърсен. В разглеждания случай тази основна задача е изпълнена успешно.

\section{Заключение}
Анализираният проект е добър пример за съвременно лично портфолио, в което визуалната идентичност и техническата реализация работят в синхрон. HTML осигурява ясна структурна рамка, CSS оформя секциите като отделни, но свързани визуални сцени, JavaScript добавя практични и емоционални интеракции, а React е ограничен до началния екран, където носи най-осезаем ефект върху възприятието.

Силата на сайта не е в прекомерна техническа сложност, а в баланса между представяне, подреденост и умерена интерактивност. Потребителят получава последователен поток: първо впечатление, професионален контекст, доказателства за развитие и ясна възможност за контакт. Този сценарий е много добре съобразен с целта на един портфолио сайт.

В обобщение може да се каже, че проектът следва актуални тенденции в web дизайна по зрял и прагматичен начин. Вместо да преследва ефектност сама по себе си, реализацията използва визуални и технически средства там, където те реално подобряват комуникацията. Именно това прави сайта убедителен както като учебен проект, така и като практическо средство за професионално представяне.

\end{document}
